\section{Limitations and honesty ledger}

The central limitation is calibration under dependence. The rank arithmetic is
strictly prior, but grid telemetry is not shown to be exchangeable. The rolling
240-window channel references adapt to local score distributions; they do not
turn the sequence into independent or identically distributed observations.
The growing outer reference also reuses history across many tests. Consequently
the empirical rank counts are descriptive, ordinary BH is nominal, and neither
coverage nor FDR is established. Block conformal methods provide one relevant
research direction \citep{chernozhukov2018dependentconformal}, but no such block
construction is implemented here. Likewise, the conditions studied for
conformal outlier testing do not transfer automatically to this online reuse
pattern \citep{bates2023conformaloutliers}.

The record fallback is transparent but weak on a long horizon. Its expected
count is an exchangeable-continuous benchmark, not a time-series guarantee
\citep{renyi1962records}, and an early extreme can suppress later records. The
zero-alert 2020 result alongside low-ranked unalerted windows empirically
exhibits this saturation. A future selector should be specified before a new
untouched period and evaluated with dependence and finite-resolution effects in
view; this paper does not retroactively substitute such a rule.

Development history also limits inference. A historical fixed-calibration
episode produced 62 alerts; about 90\% were judged seasonal-drift artifacts,
corresponding to roughly 20-fold anti-conservatism. This is a narrative
disclosure only: there is no shipped artifact for that old run, and it is not a
current comparator. The committed current comparison contains the demand-only
variant, not a reconstruction of the historical fixed-calibration episode.
Within the present algorithm, $W\in\{2,4\}$ and $K\in\{20,40\}$ were considered
using 2019 development evidence, the change from rank-PIT to robust-z was
event-informed, and the frequency stream was added after demand-only analysis.
Freezing 2020 limits subsequent adaptation but cannot erase those choices.

The relational terminology needs similar restraint. Equation
\ref{eq:relation} resembles a Gaussian copula log-density contrast, yet the
robust-z inputs are not probability-integral transformed normal scores as in a
literal copula construction \citep{sklar1959fonctions,song2000gaussiancopula}.
The score may still be useful as a fitted dependence contrast, but copula
likelihood theory does not directly calibrate it. The three channel ranks are
also dependent, so Fisher's chi-square reference is deliberately not used
\citep{fisher1932statisticalmethods}.

Sensor selection and aggregation define detectability. The frequency stream
keeps only the maximum absolute deviation within a half-hour, and all channels
are evaluated in non-overlapping two-hour windows. A short transient can lose
its timing and be dominated by sustained background extremes, whereas a broad
weather episode can produce coherent contributions across several streams.
The stream-contribution list discloses large window residuals but is descriptive
rather than a mechanistic explanation. Missing-stream robustness, alternative
aggregations, per-stream performance, ablations, synthetics, PR-AUC, and
matched-budget comparisons have not been evaluated.

The baseline suite is post-hoc and protocol-non-identical. Global scaling,
direct growing raw-score ranks, and unbudgeted record selection differ from
TRACE-C's regime conditioning, two-stage ranks, aggregation, BH attempt, and
fixed-block budget. Only TRACE-C was frozen before 2020 inspection, so apparent
wins in either direction are exploratory. Event-best ranks also favor longer
intervals through more opportunities.

The attention limit is two selected windows per fixed 48-row block. It is not a
calendar service: days with 46 or 50 settlement periods and daylight-saving
transitions are not handled as true civil days. A deployment would need a
calendar-aware budget and explicit policy for late or revised telemetry.

Finally, the governance scope is deliberately narrow. TRACE-C ranks process
and sensor anomalies in operational telemetry for human review. It is not
designed, evaluated, or authorized for person-risk scoring, and the evidence in
this study cannot support decisions about individuals.
